% GENERATED FILE. Edit canonical lessons, then run npm run generate:books.
% canonical-source-hash: fnv1a64:bdf0fd433d42b88c
% canonical-lessons: IT-C419-vagone, IT-C419-carrozza, IT-C419-scompartimento, IT-C419-cuccetta, IT-C419-capotreno

\chapter{In treno}
\label{ch:it-in-treno}
\hlchaptermodality{419}

\begin{chapteropening}
\textbf{By the end of this chapter:} \emph{I can find my place on a train: the carriage, the coach number, the compartment, the couchette and the guard.}

Complete the last lesson of chapter 419: I can find my place on a train: the carriage, the coach number, the compartment, the couchette and the guard.
\end{chapteropening}

\section[il vagone]{il vagone — the railway carriage, the wagon}

\label{lesson:IT-C419-vagone}



\begin{quote}
Before the new one: say the Italian for the ticket inspector, then the Italian for the supplement.
\end{quote}

\subsection*{You'll want to know: il vagone}
\textbf{il vagone} — \textquotedblleft{}the railway carriage, the wagon\textquotedblright{}.

\begin{quote}
Il nostro vagone è in fondo al treno.
\end{quote}

\begin{quote}
\emph{Our carriage is at the back of the train.}
\end{quote}

\emph{Il vagone letto} — sleeping car; \emph{il vagone merci} — goods wagon.

\begin{cousinweb}
\emph{vagone} was borrowed from French \emph{wagon}, which took it from English \emph{waggon}, from Dutch \emph{wagen} (EtymDB).
\end{cousinweb}

\begin{grammarlens}[title={What you've built}]
One more word for telling, explaining and describing at B1.
\end{grammarlens}

\subsection*{Your turn}
\begin{itemize}
\raggedright
  \item \emph{Say it:} \emph{il vagone}
  \item \emph{Say it:} \emph{il vagone}, once more
  \item \emph{Say it:} a sentence of your own with \emph{il vagone}
  \item \emph{Recall:} say the Italian for the ticket inspector, then the Italian for the supplement, then say \emph{il vagone} again
\end{itemize}

\begin{tcolorbox}[breakable,colback=teal!4,colframe=teal!35!black,title={Before you move on}]
What does il vagone mean? (\textquotedblleft{}The railway carriage, the wagon\textquotedblright{}.) Say it once more.
\end{tcolorbox}
\section[la carrozza]{la carrozza — the coach (of a train); the horse-drawn carriage}

\label{lesson:IT-C419-carrozza}



\begin{quote}
Before the new one: say the Italian for the supplement, then the Italian for the railway carriage.
\end{quote}

\subsection*{You'll want to know: la carrozza}
\textbf{la carrozza} — \textquotedblleft{}the coach (of a train); the horse-drawn carriage\textquotedblright{}.

\begin{quote}
Abbiamo i posti nella carrozza sette.
\end{quote}

\begin{quote}
\emph{We have seats in coach seven.}
\end{quote}

On a ticket \emph{carrozza} is the coach number; \emph{la carrozza} is also a horse-drawn carriage.

\begin{cousinweb}
\emph{carrozza} is built on \emph{carro} 'cart' (EtymDB), which is inherited from Latin \emph{carrus}, a word of Gaulish origin.
\end{cousinweb}

\begin{grammarlens}[title={What you've built}]
One more word for telling, explaining and describing at B1.
\end{grammarlens}

\subsection*{Your turn}
\begin{itemize}
\raggedright
  \item \emph{Say it:} \emph{la carrozza}
  \item \emph{Say it:} \emph{la carrozza}, once more
  \item \emph{Say it:} a sentence of your own with \emph{la carrozza}
  \item \emph{Recall:} say the Italian for the supplement, then the Italian for the railway carriage, then say \emph{la carrozza} again
\end{itemize}

\begin{tcolorbox}[breakable,colback=teal!4,colframe=teal!35!black,title={Before you move on}]
What does la carrozza mean? (\textquotedblleft{}The coach (of a train); the horse-drawn carriage\textquotedblright{}.) Say it once more.
\end{tcolorbox}
\section[lo scompartimento]{lo scompartimento — the compartment}

\label{lesson:IT-C419-scompartimento}



\begin{quote}
Before the new one: say the Italian for the railway carriage, then the Italian for the coach.
\end{quote}

\subsection*{You'll want to know: lo scompartimento}
\textbf{lo scompartimento} — \textquotedblleft{}the compartment\textquotedblright{}.

\begin{quote}
Lo scompartimento ha sei posti.
\end{quote}

\begin{quote}
\emph{The compartment has six seats.}
\end{quote}

Old trains have them; modern ones are open.

\begin{cousinweb}
\emph{scompartimento} is built on \emph{compartimento}, from Late Latin \emph{compartīrī} 'to share out' (EtymDB), a learned borrowing built on \emph{pars} 'part'.
\end{cousinweb}

\begin{grammarlens}[title={What you've built}]
One more word for telling, explaining and describing at B1.
\end{grammarlens}

\subsection*{Your turn}
\begin{itemize}
\raggedright
  \item \emph{Say it:} \emph{lo scompartimento}
  \item \emph{Say it:} \emph{lo scompartimento}, once more
  \item \emph{Say it:} a sentence of your own with \emph{lo scompartimento}
  \item \emph{Recall:} say the Italian for the railway carriage, then the Italian for the coach, then say \emph{lo scompartimento} again
\end{itemize}

\begin{tcolorbox}[breakable,colback=teal!4,colframe=teal!35!black,title={Before you move on}]
What does lo scompartimento mean? (\textquotedblleft{}The compartment\textquotedblright{}.) Say it once more.
\end{tcolorbox}
\section[la cuccetta]{la cuccetta — the couchette, the berth}

\label{lesson:IT-C419-cuccetta}



\begin{quote}
Before the new one: say the Italian for the coach, then the Italian for the compartment.
\end{quote}

\subsection*{You'll want to know: la cuccetta}
\textbf{la cuccetta} — \textquotedblleft{}the couchette, the berth\textquotedblright{}.

\begin{quote}
Di notte dormiamo in cuccetta.
\end{quote}

\begin{quote}
\emph{At night we sleep in a couchette.}
\end{quote}

On a ship too: a bunk.

\begin{cousinweb}
\emph{cuccetta} is a little \emph{cuccia} 'dog bed, bed' (EtymDB), which was borrowed from French \emph{couche} 'bed'.
\end{cousinweb}

\begin{grammarlens}[title={What you've built}]
One more word for telling, explaining and describing at B1.
\end{grammarlens}

\subsection*{Your turn}
\begin{itemize}
\raggedright
  \item \emph{Say it:} \emph{la cuccetta}
  \item \emph{Say it:} \emph{la cuccetta}, once more
  \item \emph{Say it:} a sentence of your own with \emph{la cuccetta}
  \item \emph{Recall:} say the Italian for the coach, then the Italian for the compartment, then say \emph{la cuccetta} again
\end{itemize}

\begin{tcolorbox}[breakable,colback=teal!4,colframe=teal!35!black,title={Before you move on}]
What does la cuccetta mean? (\textquotedblleft{}The couchette, the berth\textquotedblright{}.) Say it once more.
\end{tcolorbox}
\section[il capotreno]{il capotreno — the guard, the train manager}

\label{lesson:IT-C419-capotreno}



\begin{quote}
Before the new one: say the Italian for the compartment, then the Italian for the couchette.
\end{quote}

\subsection*{You'll want to know: il capotreno}
\textbf{il capotreno} — \textquotedblleft{}the guard, the train manager\textquotedblright{}.

\begin{quote}
Il capotreno annuncia la prossima fermata.
\end{quote}

\begin{quote}
\emph{The guard announces the next stop.}
\end{quote}

\emph{Capo} + \emph{treno}: the head of the train crew.

\begin{cousinweb}
\emph{capotreno} is built from \emph{capo} 'head', inherited from Latin \emph{caput}, and \emph{treno} 'train', which was borrowed from French \emph{train} (EtymDB).
\end{cousinweb}

\begin{grammarlens}[title={What you've built}]
One more word for telling, explaining and describing at B1.
\end{grammarlens}

\subsection*{Your turn}
\begin{itemize}
\raggedright
  \item \emph{Say it:} \emph{il capotreno}
  \item \emph{Say it:} \emph{il capotreno}, once more
  \item \emph{Say it:} a sentence of your own with \emph{il capotreno}
  \item \emph{Recall:} say the Italian for the compartment, then the Italian for the couchette, then say \emph{il capotreno} again
\end{itemize}

\begin{tcolorbox}[breakable,colback=teal!4,colframe=teal!35!black,title={Before you move on}]
What does il capotreno mean? (\textquotedblleft{}The guard, the train manager\textquotedblright{}.) Say it once more.
\end{tcolorbox}
